\section{Statistical Certificates and Learning Complexity}
\label{app:b}

This appendix proves the statistical results for a fixed tail level
$\alpha\in(0,1)$ and a fixed finite sensitivity level $\Gamma\geq 1$.
Write $\beta=1-\alpha$, $p=p_1$, and $P=P_{\mathrm{obs}}$; thus $Pg$
denotes integration of an observed-data function $g$ against the common
observed law. For brevity, set
\[
 \rho^L(f)=\underline{\rho}_{\alpha,\Gamma}(f),\qquad
 \rho^U(f)=\overline{\rho}_{\alpha,\Gamma}(f),\qquad
 I_\Gamma(f)=\rho^U(f)-\rho^L(f).
\]
The population characterization in Appendix~\ref{app:a} gives
$\rho^\sigma(f)=\inf_{\eta\in[0,1]}
\{\eta+\beta^{-1}m^\sigma_{f,\eta}\}$ for
$\sigma\in\{L,U\}$. All $L_2$ norms below are with respect to $P_X$.

For uniform empirical-process statements, assume the paired maps
$x\mapsto(\ell(f(x),0),\ell(f(x),1))$, $f\in\mathcal F$, form a pointwise-separable
class of measurable functions with the envelope supplied by
Assumption~\ref{ass:regularity}. A countable pointwise-dense subclass and
rational thresholds including $0,1$ determine the suprema below, since the
score transformations are continuous. Measurable, pointwise-separable
versions of the fitted score classes suffice as well.

\subsection{Score regularity and exact nuisance-bias expansion}
\label{app:score-bias}

Recall $a^y_{f,\eta}(x)=(\ell(f(x),y)-\eta)_+$ and
$d_{f,\eta}(x)=a^1_{f,\eta}(x)-a^0_{f,\eta}(x)$. For a real number $d$,
write $d_+=\max(d,0)$ and $d_-=\max(-d,0)$. Define
\begin{align}
 H_p^U(d)&=q_+(p)d_+-q_-(p)d_-,&
 H_p^L(d)&=q_-(p)d_+-q_+(p)d_-, \label{eq:b:H}\\
 G_p^U(d)&=q_+'(p)d_+-q_-'(p)d_-,&
 G_p^L(d)&=q_-'(p)d_+-q_+'(p)d_- . \label{eq:b:G}
\end{align}
Consequently,
$b^\sigma_{f,\eta}(x;p)=a^0_{f,\eta}(x)+H_p^\sigma(d_{f,\eta}(x))$
and $\dot b^\sigma_{f,\eta}(x;p)=G_p^\sigma(d_{f,\eta}(x))$.

\begin{lemma}[Regularity of the endpoint maps]
\label{lemma:b:regularity}
For $p\in[0,1]$, the derivatives of the sensitivity endpoints satisfy
\begin{align}
 q_+'(p)&=\frac{\Gamma}{\{1+(\Gamma-1)p\}^2},&
 q_-'(p)&=\frac{\Gamma}{\{\Gamma-(\Gamma-1)p\}^2},\label{eq:b:first-derivatives}\\
 q_+''(p)&=-\frac{2\Gamma(\Gamma-1)}{\{1+(\Gamma-1)p\}^3},&
 q_-''(p)&=\frac{2\Gamma(\Gamma-1)}{\{\Gamma-(\Gamma-1)p\}^3}.
 \label{eq:b:second-derivatives}
\end{align}
In particular, $\Gamma^{-1}\leq q_\pm'(p)\leq\Gamma$ and
$|q_\pm''(p)|\leq 2\Gamma(\Gamma-1)$. Uniformly in $(f,\eta,x,p,\sigma)$,
\begin{equation}
 \begin{gathered}
 0\leq b^\sigma_{f,\eta}(x;p)\leq1,\qquad
 |\dot b^\sigma_{f,\eta}(x;p)|\leq\Gamma,\\
 |\partial_p^2 b^\sigma_{f,\eta}(x;p)|\leq2\Gamma(\Gamma-1).
 \end{gathered}
 \label{eq:b:regularity-bounds}
\end{equation}
Moreover, $H_p^\sigma$ is $1$-Lipschitz and $G_p^\sigma$ is
$\Gamma$-Lipschitz on the real line; both vanish at zero.
\end{lemma}

\begin{proof}
Differentiating $q_+(p)=\Gamma p/\{1+(\Gamma-1)p\}$ and
$q_-(p)=p/\{\Gamma-(\Gamma-1)p\}$ gives
\eqref{eq:b:first-derivatives}--\eqref{eq:b:second-derivatives}.
The denominators before taking powers lie in $[1,\Gamma]$, giving the
derivative bounds. The value $a^0+H_p^\sigma(a^1-a^0)$ is a convex
combination of $a^0,a^1\in[0,1]$; its derivatives are bounded by
$\Gamma|a^1-a^0|$ and $2\Gamma(\Gamma-1)|a^1-a^0|$, respectively.
This proves \eqref{eq:b:regularity-bounds}. On either half-line the slopes
of $H_p^\sigma$ and $G_p^\sigma$ lie in $[0,1]$ and
$[\Gamma^{-1},\Gamma]$, respectively. Both maps are continuous at zero,
so these slope bounds also prove the Lipschitz bounds across loss ties.
\end{proof}

For candidate nuisance functions $p_c:\mathcal X\to[0,1]$ and
$\pi_c:\mathcal X\to[\pi_{\min}/2,1]$, the score is
\begin{align}
 \phi^\sigma_{f,\eta}(O;p_c,\pi_c)
 ={}&S a^{Y^{\mathrm{obs}}}_{f,\eta}(X)
 +(1-S)b^\sigma_{f,\eta}(X;p_c(X)) \notag\\
 &+S\frac{1-\pi_c(X)}{\pi_c(X)}
 \dot b^\sigma_{f,\eta}(X;p_c(X))
 \{Y^{\mathrm{obs}}-p_c(X)\}.
 \label{eq:b:score}
\end{align}
The sum of its first two terms belongs to $[0,1]$, whereas its correction
has absolute value at most $2\Gamma/\pi_{\min}$. Therefore
\begin{equation}
 \begin{split}
 -\frac{2\Gamma}{\pi_{\min}}&\leq
 \phi^\sigma_{f,\eta}(O;p_c,\pi_c)
 \leq1+\frac{2\Gamma}{\pi_{\min}},\\
 |\phi^\sigma_{f,\eta}(O;p_c,\pi_c)|&\leq
 B:=1+\frac{2\Gamma}{\pi_{\min}}.
 \end{split}
 \label{eq:b:score-envelope}
\end{equation}
The score can be negative; the concentration argument below explicitly
allows signed classes.

\begin{lemma}[Exact conditional nuisance bias]
\label{lemma:b:bias}
Suppress $(f,\eta,\sigma,X)$ inside $b$ and its derivatives. With
$e=p_c-p$, the conditional difference between the score expectation and
the target excess-loss integrand is exactly
\begin{equation}
 (1-\pi)\{b(p_c)-b(p)-\dot b(p_c)e\}
 +\frac{\pi_c-\pi}{\pi_c}\dot b(p_c)e.
 \label{eq:b:conditional-bias}
\end{equation}
Taking $P_X$ expectations of \eqref{eq:b:conditional-bias} gives
$P\phi^\sigma_{f,\eta}(\cdot;p_c,\pi_c)-m^\sigma_{f,\eta}$, and
\begin{equation}
 \begin{aligned}
 &\left|P\phi^\sigma_{f,\eta}(\cdot;p_c,\pi_c)-m^\sigma_{f,\eta}\right|\\
 &\quad\leq\Gamma(\Gamma-1)\|p_c-p\|_2^2
 +\frac{2\Gamma}{\pi_{\min}}\|p_c-p\|_2\|\pi_c-\pi\|_2.
 \end{aligned}
 \label{eq:b:bias-bound}
\end{equation}
\end{lemma}

\begin{proof}
Conditional on $X$, the observed excess-loss contribution has mean
$\pi\{(1-p)a^0+pa^1\}$ and
$\mathbb E[S(Y^{\mathrm{obs}}-p_c)\mid X]=\pi(p-p_c)$.
Subtracting the target integrand
$\pi\{(1-p)a^0+pa^1\}+(1-\pi)b(p)$ leaves
\[
 (1-\pi)\{b(p_c)-b(p)\}
 +\frac{\pi(1-\pi_c)}{\pi_c}\dot b(p_c)(p-p_c).
\]
The identity
$\pi(1-\pi_c)/\pi_c=(1-\pi)+(\pi-\pi_c)/\pi_c$
reduces this expression to \eqref{eq:b:conditional-bias}.
Taylor's theorem, expanded about $p_c$ along the segment joining $p_c$
and $p$, and Lemma~\ref{lemma:b:regularity} give
\[
 |b(p_c)-b(p)-\dot b(p_c)(p_c-p)|
 \leq\Gamma(\Gamma-1)(p_c-p)^2.
\]
Use $0\leq1-\pi\leq1$ for the first term, and
$|\dot b(p_c)/\pi_c|\leq2\Gamma/\pi_{\min}$ followed by
Cauchy--Schwarz for the second term. This yields \eqref{eq:b:bias-bound}
uniformly over all score indices.
\end{proof}

The score is unbiased at the true nuisances and has $O(t^2)$ bias along
admissible nuisance paths with $L_2$ perturbations of order $t$.
The sign of $d_{f,\eta}(x)$ is independent of the nuisance argument, so
endpoint selection adds no remainder. At $\Gamma=1$ the curvature term
vanishes; for $\Gamma>1$, knowing $\pi$ alone generally leaves quadratic
error in $p$. Thus the result is an orthogonal expansion, not generic
double robustness. Its overlap premise and clipping ranges are substantive:
clipping cannot supply labels or identification where overlap fails.

\subsection{Cross-fitting and uniform empirical-process control}
\label{app:empirical-process}

Let $I_1,\ldots,I_K$ be a deterministic partition of the observation indices,
with $n_k=|I_k|\geq1$, $n_{\min}=\min_k n_k$, and $w_k=n_k/n$.
A random partition independent of the observations is also allowed by
conditioning on that partition. Write $D_{-k}$ for the training observations
outside $I_k$ and $P_{n,k}$ for the empirical measure on $I_k$.
The clipped fits $\widehat p^{(-k)},\widehat\pi^{(-k)}$ are measurable
functions of $D_{-k}$. Any independent fitting randomization can be included
in this conditioning. Define
\[
 \phi^{\sigma,(k)}_{f,\eta}
 =\phi^\sigma_{f,\eta}
 (\cdot;\widehat p^{(-k)},\widehat\pi^{(-k)}),\qquad
 \Phi_k^\sigma=\{\phi^{\sigma,(k)}_{f,\eta}:f\in\mathcal F,
 \ \eta\in[0,1]\}.
\]
The empirical Rademacher complexities used here are
\begin{equation}
 \begin{split}
 \widehat{\mathfrak R}^{\sigma,k}_{\mathrm{abs}}
 &=\mathbb E_\varepsilon\!\left[
 \sup_{f,\eta}\left|\frac1{n_k}
 \sum_{i\in I_k}\varepsilon_i\phi^{\sigma,(k)}_{f,\eta}(O_i)\right|
 \,\middle|\,D_n\right],\\
 \widehat{\mathfrak R}_{n,\mathrm{abs}}^\sigma
 &=\sum_{k=1}^K w_k\widehat{\mathfrak R}^{\sigma,k}_{\mathrm{abs}}.
 \end{split}
 \label{eq:b:empirical-rademacher}
\end{equation}
Here $D_n$ includes the observations, partition, and fitting seeds.
The expectation averages only independent uniform signs in $\{-1,1\}$;
the fitted functions and evaluation observations are fixed.

\begin{lemma}[Empirical-complexity bound for a signed class]
\label{lemma:b:empirical-rademacher}
Let $Z_1,\ldots,Z_m$ be independent with law $Q$ and empirical measure $Q_m$,
and let $\mathcal H$ be a fixed pointwise-separable class with $|h|\leq B$. For $0<\xi<1$,
with probability at least $1-\xi$,
\begin{equation}
 \sup_{h\in\mathcal H}|(Q_m-Q)h|
 \leq2\widehat{\mathfrak R}_{m,\mathrm{abs}}(\mathcal H)
 +6B\sqrt{\frac{\log(2/\xi)}{2m}},
 \label{eq:b:signed-concentration}
\end{equation}
where the empirical complexity uses the absolute supremum as in
\eqref{eq:b:empirical-rademacher}.
\end{lemma}

\begin{proof}
Put $Z=\sup_h|(Q_m-Q)h|$ and
$R=\widehat{\mathfrak R}_{m,\mathrm{abs}}(\mathcal H)$.
Symmetrization gives $\mathbb EZ\leq2\mathbb ER$
\cite{bartlett2002rademacher,van1996weak}. Changing one observation changes
each of $Z,R$ by at most $2B/m$. With
$t=2B\sqrt{\log(2/\xi)/(2m)}$, McDiarmid's inequality bounds each of
$\mathbb P(Z>\mathbb EZ+t)$ and $\mathbb P(\mathbb ER>R+t)$ by $\xi/2$.
Thus $Z\leq2R+3t$ with probability at least $1-\xi$, proving
\eqref{eq:b:signed-concentration} for signed classes with $|h|\leq B$.
\end{proof}

Conditional on $D_{-k}$, the observations in $I_k$ are independent with
law $P$, and $\Phi_k^\sigma$ is fixed with envelope $B$.
Apply Lemma~\ref{lemma:b:empirical-rademacher} with
$m=n_k$ and $\xi=\delta/(2K)$. Its conditional failure probability is
at most $\delta/(2K)$, so the same bound holds unconditionally.
A union bound over the $2K$ fold--endpoint pairs yields an event
$\mathcal E_{\mathrm{emp}}$ of probability at least $1-\delta$ on which
\begin{align}
 \sup_{f,\eta}\left|
 \sum_{k=1}^K w_k(P_{n,k}-P)\phi^{\sigma,(k)}_{f,\eta}
 \right|
 &\leq2\widehat{\mathfrak R}_{n,\mathrm{abs}}^\sigma
 +6B\sum_{k=1}^K w_k\sqrt{\frac{\log(4K/\delta)}{2n_k}}\notag\\
 &\leq2\widehat{\mathfrak R}_{n,\mathrm{abs}}^\sigma
 +6B\sqrt{\frac{\log(4K/\delta)}{2n_{\min}}}
 \label{eq:b:fold-bound}
\end{align}
for both endpoints. Here $P\phi^{\sigma,(k)}_{f,\eta}$ integrates the
fitted function against a fresh observation. The training sets overlap;
no independence between folds or between their success events is used.
In particular, the proof never conditions simultaneously on all fitted
nuisances, which could destroy the needed held-out independence.

\subsection{Uniform endpoint estimation and complexity rates}
\label{app:uniform-estimation}

Define the cross-fitted empirical objectives and their infima by
\begin{equation}
 \begin{aligned}
 \widehat Q_f^\sigma(\eta)&=\eta+\frac1\beta
 \sum_{k=1}^K w_kP_{n,k}\phi^{\sigma,(k)}_{f,\eta},\\
 \widehat\rho^\sigma(f)&=\inf_{\eta\in[0,1]}\widehat Q_f^\sigma(\eta).
 \end{aligned}
 \label{eq:b:empirical-objectives}
\end{equation}

\begin{proof}[Proof of Theorem~\ref{thm:learning}]
Under the theorem's assumptions and the measurability convention above,
let $\mathcal E_\nu$ be the simultaneous nuisance event
\begin{equation}
 \max_k\|\widehat p^{(-k)}-p\|_2\leq r_p,\qquad
 \max_k\|\widehat\pi^{(-k)}-\pi\|_2\leq r_\pi,
 \label{eq:b:nuisance-event}
\end{equation}
with probability at least $1-\delta_\nu$. For
$\delta\in(0,1)$ and $\delta_\nu\in[0,1-\delta)$, the
explicit calibrated radius is
\begin{equation}
 \begin{aligned}
 r_n(\delta)=\frac1\beta\Bigg[
 &2\max_{\sigma\in\{L,U\}}\widehat{\mathfrak R}_{n,\mathrm{abs}}^\sigma
 +6B\sqrt{\frac{\log(4K/\delta)}{2n_{\min}}}\\
 &+\Gamma(\Gamma-1)r_p^2
 +\frac{2\Gamma}{\pi_{\min}}r_pr_\pi\Bigg].
 \end{aligned}
 \label{eq:b:radius}
\end{equation}
On $\mathcal E_\nu$, Lemma~\ref{lemma:b:bias}, applied to each realized
pair of clipped fits, implies
\[
 \sup_{f,\eta,\sigma}\left|
 \sum_{k=1}^K w_kP\phi^{\sigma,(k)}_{f,\eta}
 -m^\sigma_{f,\eta}\right|
 \leq\Gamma(\Gamma-1)r_p^2+\frac{2\Gamma}{\pi_{\min}}r_pr_\pi.
\]
The intersection $\mathcal E=\mathcal E_\nu\cap\mathcal E_{\mathrm{emp}}$
has probability at least $1-\delta-\delta_\nu$ by the union bound.
On this event, \eqref{eq:b:fold-bound} and the last display give
\[
 \sup_{f,\eta,\sigma}\left|
 \sum_{k=1}^K w_kP_{n,k}\phi^{\sigma,(k)}_{f,\eta}
 -m^\sigma_{f,\eta}\right|\leq\beta r_n(\delta).
\]
Applying the infimum stability inequality \eqref{eq:a:infimum-stability}
to $\widehat Q_f^\sigma$ and its population counterpart yields, on the same
event of probability at least $1-\delta-\delta_\nu$,
\begin{equation}
 \sup_{f\in\mathcal F}\max_{\sigma\in\{L,U\}}
 |\widehat\rho^\sigma(f)-\rho^\sigma(f)|\leq r_n(\delta).
 \label{eq:b:uniform-endpoints}
\end{equation}
This requires neither convexity of the signed empirical objective nor a
unique threshold or loss density there.
\end{proof}

For a total confidence level $1-\delta_{\mathrm{tot}}$, allocate
$\delta+\delta_\nu\leq\delta_{\mathrm{tot}}$. The empirical complexity in
\eqref{eq:b:radius} is an expectation over all Rademacher signs, conditional
on the realized data. A Monte Carlo estimate is not automatically an upper
bound for this expectation. Any random complexity or nuisance-error bound
used instead must hold simultaneously on a specified validity event, whose
failure probability is included in the confidence allocation.

We next bound the score complexity. For a sample $x_1,\ldots,x_m$,
let $\widehat{\mathfrak R}_{m,\mathrm{abs}}(\mathcal V)$ denote its
absolute empirical complexity. Define
\[
 \mathcal L_y=\{x\mapsto\ell(f(x),y):f\in\mathcal F\},\qquad
 \mathcal A_y=\{x\mapsto(\ell(f(x),y)-\eta)_+:
 f\in\mathcal F,\ \eta\in[0,1]\}.
\]
Define the threshold-augmented loss-class complexity by
\begin{equation}
 \mathfrak C_n=\sum_{k=1}^K w_k\left\{
 \widehat{\mathfrak R}_{n_k,\mathrm{abs}}(\mathcal L_0)
 +\widehat{\mathfrak R}_{n_k,\mathrm{abs}}(\mathcal L_1)
 +n_k^{-1/2}\right\},
 \label{eq:b:loss-complexity}
\end{equation}
where each term uses the covariates in its evaluation fold.
For fixed fitted nuisances and $o=(x,s,y^{\mathrm{obs}})$, put
$p_c=\widehat p^{(-k)}(x)$ and $w=(1-\widehat\pi^{(-k)}(x)) /
\widehat\pi^{(-k)}(x)$. The score is
$\Psi_o^\sigma(a^0_{f,\eta}(x),a^1_{f,\eta}(x))$, where
\begin{align}
 \Psi_o^\sigma(u_0,u_1)
 ={}&s\{(1-y^{\mathrm{obs}})u_0+y^{\mathrm{obs}}u_1\}\notag\\
 &+(1-s)\{u_0+H_{p_c}^\sigma(u_1-u_0)\}\notag\\
 &+swG_{p_c}^\sigma(u_1-u_0)(y^{\mathrm{obs}}-p_c).
 \label{eq:b:composition}
\end{align}
It vanishes at $(0,0)$. Lemma~\ref{lemma:b:regularity},
$w\leq2/\pi_{\min}$, and
$|\Delta u_1-\Delta u_0|\leq2\|\Delta u\|_\infty$ show that it is
$L_*=3+4\Gamma/\pi_{\min}$-Lipschitz in the maximum norm,
and hence also in Euclidean norm.

The vector-contraction inequality \cite{maurer2016vector} therefore gives,
for a universal constant $C$,
\begin{equation}
 \begin{aligned}
 \widehat{\mathfrak R}^{\sigma,k}_{\mathrm{abs}}
 \leq C L_*\Big\{&
 \widehat{\mathfrak R}_{n_k,\mathrm{abs}}(\mathcal L_0)
 +\widehat{\mathfrak R}_{n_k,\mathrm{abs}}(\mathcal L_1)\\
 &+n_k^{-1/2}\Big\}.
 \end{aligned}
 \label{eq:b:score-complexity}
\end{equation}
To handle the absolute supremum, use that the score class contains zero
at $\eta=1$: the absolute supremum is at most the sum of the positive and
negative suprema, whose expectations agree by sign symmetry. Vector
contraction on the positive supremum, followed by separating the two
coordinate suprema, bounds it by a universal multiple of
$L_*\sum_y\widehat{\mathfrak R}_{n_k,\mathrm{abs}}(\mathcal A_y)$.
Apply the same argument to the $\sqrt2$-Lipschitz hinge map
$(z,t)\mapsto(z-t)_+$, adjoining zero to its input class if needed.
Its coordinate complexities are those of $\mathcal L_y$ and the constant
threshold class; the latter is
$m^{-1}\mathbb E_\varepsilon|\sum_i\varepsilon_i|\leq m^{-1/2}$.
This proves \eqref{eq:b:score-complexity}. All contraction maps are
continuous endpoint values or derivatives, including at ties; no
discontinuous selector class is introduced.

If $\mathcal L_y$ has VC-subgraph dimension at most $d\geq1$, then
$\log N(\epsilon,\mathcal L_y,L_2(Q))\leq C_0d\log(C_1/\epsilon)$
for $0<\epsilon<1$ and every probability measure $Q$.
After adjoining zero and the negative class, the standard entropy-to-complexity
bound \cite{van1996weak,bartlett2002rademacher} gives, uniformly over samples,
\begin{equation}
 \widehat{\mathfrak R}_{m,\mathrm{abs}}(\mathcal L_y)
 \leq C'\sqrt{\frac{d+1}{m}}.
 \label{eq:b:VC-complexity}
\end{equation}
It therefore applies to the fixed fitted nuisances in each conditional argument.

For fixed $K$ and balanced folds, $n_{\min}\geq c_Kn$ for a positive
constant $c_K$ once all folds are nonempty. Equations
\eqref{eq:b:score-complexity}--\eqref{eq:b:VC-complexity} give
\begin{equation}
 \begin{aligned}
 &\frac1\beta\left[
 2\max_\sigma\widehat{\mathfrak R}_{n,\mathrm{abs}}^\sigma
 +6B\sqrt{\frac{\log(4K/\delta)}{2n_{\min}}}\right]\\
 &\qquad\leq \frac{C_{\Gamma,\pi_{\min},K}}{\beta}
 \sqrt{\frac{d+\log(1/\delta)}{n}}.
 \end{aligned}
 \label{eq:b:VC-rate}
\end{equation}
For fixed $\Gamma,\pi_{\min},\beta$, nuisance rates $o(n^{-1/4})$ make
the nuisance contribution in \eqref{eq:b:radius} $o(n^{-1/2})$.
Joint high-probability nuisance bounds remain necessary; cross-fitting
does not establish them. All uniformity is over $f,\eta$ at fixed
$\alpha,\Gamma$.

\begin{example}[Conditional sparse-nuisance specialization]
\label{ex:b:sparse-nuisances}
Suppose sparse logistic nuisance fits in $D$ covariates satisfy simultaneous
$L_2(P_X)$ error bounds of orders
\[
 r_p\lesssim\sqrt{\frac{s_p\log(eD)}{n_{\mathrm{sel}}}},\qquad
 r_\pi\lesssim\sqrt{\frac{s_\pi\log(eD)}{n_{\mathrm{tr}}}},
\]
where $s_p,s_\pi$ are sparsities, $n_{\mathrm{tr}}$ is the smallest training-fold
size, and $n_{\mathrm{sel}}>0$ is a simultaneous lower bound for the number of
revealed training labels in every fold. Substitution in the proved remainder
gives, up to the constants in these assumed rates,
\[
 \frac1\beta\left\{
 \Gamma(\Gamma-1)\frac{s_p\log(eD)}{n_{\mathrm{sel}}}
 +\frac{2\Gamma}{\pi_{\min}}
 \frac{\sqrt{s_ps_\pi}\log(eD)}{\sqrt{n_{\mathrm{sel}}n_{\mathrm{tr}}}}
 \right\}.
\]
For a fixed number of balanced folds with nonempty training sets,
$n_{\mathrm{tr}}\asymp n$. These substitutions
require the stated model, design, tuning, and confidence conditions.
A selected-population bound transfers to $P_X$ because, for every measurable $e$,
\begin{equation}
 \|e\|_{L_2(P_X)}^2
 \leq\frac{P(S=1)}{\pi_{\min}}
 \|e\|_{L_2(P_{X\mid S=1})}^2.
 \label{eq:b:selected-norm-transfer}
\end{equation}
Indeed, $\mathbb E[\pi(X)e(X)^2]
=P(S=1)\mathbb E[e(X)^2\mid S=1]$ and $\pi(X)\geq\pi_{\min}$.
Learning $p$ from revealed outcomes must therefore control target-population error.
\end{example}

\begin{proposition}[Finite threshold search for the fitted objectives]
\label{prop:b:threshold-search}
For a fixed $f$ and realized nuisance fits, each objective in
\eqref{eq:b:empirical-objectives} attains its minimum on
\[
 T_n(f)=\{0,1\}\cup
 \{\ell(f(X_i),y):1\leq i\leq n,\ y\in\{0,1\}\}.
\]
\end{proposition}

\begin{proof}
Thresholding preserves the label-loss ordering at each $X_i$.
Thus $H_{p_c}^\sigma(d_{f,\eta})$ and $G_{p_c}^\sigma(d_{f,\eta})$
are fixed coefficients times $d_{f,\eta}$ throughout $[0,1]$; at
$d_{f,\eta}=0$ both representations vanish. Each score is a fixed linear
combination of the two label-specific hinges, making the objective continuous
and affine between successive knots of $T_n(f)$. Each such interval has
an endpoint minimizer, even if the objective is constant there.
Signed correction coefficients may prevent convexity but preserve this result.
\end{proof}

An exact implementation can sort the at most $2n+2$ breakpoints once
per predictor. Write the fitted score as
$c_{i0}^\sigma(\ell(f(X_i),0)-\eta)_+
+c_{i1}^\sigma(\ell(f(X_i),1)-\eta)_+$, with coefficients fixed across
$\eta$. Starting from the objective at zero, sweep the sorted breakpoints:
the objective is updated by its current slope times the distance to the
next breakpoint, and the slope increases by $c_{iy}^\sigma/(n\beta)$
when the corresponding hinge vanishes. Combine all coefficients at a
tied breakpoint and initialize the right slope using only losses above
zero. Checking all knot values gives both global minima in
$O(n\log n)$ time and $O(n)$ storage, after the fitted probabilities
and label-specific losses have been evaluated.

\subsection{Outer certificates, comparisons, and predictor selection}
\label{app:decision-guarantees}

\begin{proof}[Proof of the decision guarantees in Section~\ref{sec:uniform-certificates}]
Work on the joint event \eqref{eq:b:uniform-endpoints}, of probability at
least $1-\delta-\delta_\nu$, and use the certificate
\eqref{eq:outer-certificate}. Since $0\leq\rho^L(f)\leq\rho^U(f)\leq1$,
clipping preserves containment
of the sharp interval and hence every compatible CVaR. Containment gives
the lower width bound; clipping decreases width, while
\[
 \widehat\rho^U(f)-\widehat\rho^L(f)+2r_n
 \leq\rho^U(f)-\rho^L(f)+4r_n,
\]
proving the upper bound. If
$\widehat\rho^U(f)+r_n<\widehat\rho^L(g)-r_n$, then for every compatible law,
\begin{align*}
 \rho_{\alpha,P^\star}(f)&\leq\rho^U(f)
 \leq\widehat\rho^U(f)+r_n\\
 &<\widehat\rho^L(g)-r_n\leq\rho^L(g)
 \leq\rho_{\alpha,P^\star}(g).
\end{align*}
Finally, any measurable $\varepsilon_{\mathrm{opt}}$-approximate minimizer
$\widehat f\in\mathcal F$ of $\widehat\rho^U$ satisfies
\begin{align*}
 \rho^U(\widehat f)&\leq\inf_f\widehat\rho^U(f)+\varepsilon_{\mathrm{opt}}+r_n\\
 &\leq\inf_f\rho^U(f)+\varepsilon_{\mathrm{opt}}+2r_n.
\end{align*}
The last inequality takes infima of the pointwise endpoint bound, so no
exact minimizer is needed. Uniformity permits all predictors in these
claims to be chosen using the same data.
\end{proof}

The selection benchmark is the smallest compatible upper endpoint in
the class. For every compatible $P^\star$, the same argument gives
$\rho_{\alpha,P^\star}(\widehat f)\leq
\inf_f\rho^U(f)+2r_n+\varepsilon_{\mathrm{opt}}$.
Clipped endpoint errors are at most $2r_n$ on the same event. Convergence
in probability follows if $r_n\to0$ in probability and
$\delta_n+\delta_{\nu,n}\to0$; a fixed failure probability remains fixed.

\subsection{Oracle calibration for finite-cell experiments}
\label{app:oracle-calibration}

For the oracle coverage experiment, suppose $p=p_1$ is known and define
\begin{equation}
 \begin{split}
 h^\sigma_{f,\eta}(O)&=S a^{Y^{\mathrm{obs}}}_{f,\eta}(X)
 +(1-S)b^\sigma_{f,\eta}(X;p(X)),\\
 \widehat\rho^\sigma_{\mathrm{or}}(f)&=
 \inf_{\eta\in[0,1]}\{\eta+\beta^{-1}P_nh^\sigma_{f,\eta}\}.
 \end{split}
 \label{eq:b:oracle-score}
\end{equation}
This score lies in $[0,1]$ and has mean $m^\sigma_{f,\eta}$; estimating
$\pi$ is unnecessary for this oracle construction.

\begin{proposition}[Finite-cell oracle certificate]
\label{prop:b:oracle-calibration}
Fix a predictor $f$ independently of the evaluation observations.
Suppose $P_X$ is supported on $J$ covariate values $x_1,\ldots,x_J$.
Let
\[
 T(f)=\{0,1\}\cup\{\ell(f(x_j),y):1\leq j\leq J,
 \ y\in\{0,1\}\},\qquad M=|T(f)|\leq2J+2.
\]
With probability at least $1-\delta$,
\begin{equation}
 \max_{\sigma\in\{L,U\}}
 |\widehat\rho^\sigma_{\mathrm{or}}(f)-\rho^\sigma(f)|
 \leq r_{\mathrm{or}}(\delta):=
 \frac1\beta\sqrt{\frac{\log(4M/\delta)}{2n}}.
 \label{eq:b:oracle-radius}
\end{equation}
Thus the clipped interval obtained from \eqref{eq:outer-certificate}
using these oracle estimators and $r_{\mathrm{or}}$ contains the entire
sharp population interval with probability at least $1-\delta$.
\end{proposition}

\begin{proof}
By the piecewise-affine argument of Proposition~\ref{prop:b:threshold-search},
both $h^\sigma_{f,\eta}(O)$ and its expectation have knots only in $T(f)$.
The absolute empirical-minus-population difference is maximal at a knot:
\[
 \sup_{\eta\in[0,1]}|(P_n-P)h^\sigma_{f,\eta}|
 =\max_{\eta\in T(f)}|(P_n-P)h^\sigma_{f,\eta}|.
\]
Hoeffding's inequality and a union bound over the $2M$ endpoint--threshold
pairs bound the failure probability at tolerance $u$ by $4Me^{-2nu^2}$.
Taking $u=\sqrt{\log(4M/\delta)/(2n)}$ and applying
\eqref{eq:a:infimum-stability} proves \eqref{eq:b:oracle-radius};
The containment argument in Appendix~\ref{app:decision-guarantees} gives
whole-interval coverage.
\end{proof}

The population-support set $T(f)$ includes losses at cells absent from a
particular sample; observed knots $T_n(f)$ suffice for empirical minimization
but cannot replace $T(f)$ in this concentration argument.
Conditioning covers a predictor trained on independent data. Selection using
evaluation data requires a uniform bound: for a fixed finite family
$\mathcal F_0$, replace $M$ by $\sum_{f\in\mathcal F_0}|T(f)|$.
Known $p$ supplies the oracle experiment's premise; coverage for learned-nuisance
estimators instead requires Theorem~\ref{thm:learning}'s controls.

\subsection{Revealed-tail learning lower bound}
\label{app:learning-lower-bound}

Equal tail-cell masses and constant observation probability specify the lower-bound
construction; the general model allows covariate-dependent observation and
heterogeneous revelation costs.

\begin{theorem}[Revealed-tail observation complexity]
\label{thm:main:lower-bound}
Fix $\alpha\in(0,1)$, $\beta=1-\alpha$, $\pi_0\in(0,1]$, and $n\geq1$.
For every integer $1\leq d\leq n\pi_0\beta$, there exist a binary
classification class $\mathcal F$ of VC dimension $d$ and a family
$\mathcal P$ of full-data laws with $\Gamma=1$ and $\pi(X)=\pi_0$ such that
\begin{equation}
 \begin{aligned}
 &\inf_{\widehat f}\sup_{P^\star\in\mathcal P}
 \mathbb E_{(P^\star_{\mathrm{obs}})^{\otimes n}}
 \left[\rho_{\alpha,P^\star}(\widehat f)
       -\inf_{f\in\mathcal F}\rho_{\alpha,P^\star}(f)\right]\\
 &\qquad\geq\frac1{16}\sqrt{\frac{d}{n\pi_0\beta}}.
 \end{aligned}
 \label{eq:b:learning-lower-bound}
\end{equation}
The infimum ranges over measurable learners taking values in $\mathcal F$,
including randomized learners.
\end{theorem}

\begin{proof}
We use a direct binary-testing form of the hypercube argument
\cite{yu1997assouad}. Fix
$\pi_0\in(0,1]$, $n\geq1$, and an integer
$1\leq d\leq n\pi_0\beta$.
Put
\begin{equation}
 \gamma=\frac18\sqrt{\frac{d}{n\pi_0\beta}}\leq\frac18.
 \label{eq:b:lower-gamma}
\end{equation}
Let $\mathcal X=\{\star,1,\ldots,d\}$ with
$P_X(\star)=1-\beta$ and $P_X(j)=\beta/d$.
For $v\in\{-1,1\}^d$, define
$f_v(\star)=0$ and $f_v(j)=(1+v_j)/2$, and set
$\mathcal F=\{f_v:v\in\{-1,1\}^d\}$.
The class shatters $\{1,\ldots,d\}$, and cannot shatter any larger set
because its value at $\star$ is fixed. Its VC dimension is exactly $d$.
Use ordinary zero-one loss $\ell(f(x),y)=\mathbf1\{f(x)\ne y\}$.

For each $\theta\in\{-1,1\}^d$, construct a full-data law $P_\theta$
with $Y=0$ at $\star$ and
\[
 P_\theta(Y=1\mid X=j)=\frac12+\gamma\theta_j,
 \qquad j=1,\ldots,d.
\]
Let $S\sim\operatorname{Bernoulli}(\pi_0)$ independently of $(X,Y)$,
and observe $(X,S,SY)$. These laws satisfy overlap with
$\pi(X)=\pi_0$ and conditional ignorability, hence belong to the
$\Gamma=1$ model.
All losses vanish at $\star$. The conditional error within
$\{1,\ldots,d\}$ is
\begin{equation}
 R_\theta(f_v)=\frac12-\gamma+
 \frac{2\gamma}{d}\sum_{j=1}^d\mathbf1\{v_j\ne\theta_j\}.
 \label{eq:b:hypercube-risk}
\end{equation}
The full-population loss is Bernoulli with parameter
$q_{\theta,v}=\beta R_\theta(f_v)\leq\beta$; the Bernoulli identity
\eqref{eq:a:bernoulli-cvar} gives
\begin{equation}
 \begin{split}
 \rho_{\alpha,P_\theta}(f_v)&=R_\theta(f_v),\\
 \rho_{\alpha,P_\theta}(f_v)-\inf_{f\in\mathcal F}
 \rho_{\alpha,P_\theta}(f)
 &=\frac{2\gamma}{d}\sum_{j=1}^d\mathbf1\{v_j\ne\theta_j\}.
 \end{split}
 \label{eq:b:CVaR-hamming}
\end{equation}
The minimizing predictor is $f_\theta$.

Let $Q_\theta=P_{\theta,\mathrm{obs}}^{\otimes n}$ and let
$\theta^{(j)}$ be obtained by reversing the $j$th coordinate.
For one observation, the two distributions differ only on
$\{X=j,S=1\}$, which has probability $\beta\pi_0/d$ under both laws.
Conditioning on $(X,S)$ and using the chain rule for relative entropy
therefore gives
\begin{align}
 D_{\mathrm{KL}}(Q_\theta\Vert Q_{\theta^{(j)}})
 &=\frac{n\beta\pi_0}{d}
 D_{\mathrm{KL}}\!\left(
 \operatorname{Bernoulli}(1/2+\gamma)
 \,\Vert\,\operatorname{Bernoulli}(1/2-\gamma)\right)\notag\\
 &\leq\frac{16n\beta\pi_0\gamma^2}{d}=\frac14.
 \label{eq:b:KL}
\end{align}
Here the symmetric Bernoulli divergences both equal
$2\gamma\log\{(1+2\gamma)/(1-2\gamma)\}\leq16\gamma^2$ for
$\gamma\leq1/4$, and independence supplies the factor $n$.
Pinsker's inequality gives
\begin{equation}
 \operatorname{TV}(Q_\theta,Q_{\theta^{(j)}})
 \leq\sqrt{D_{\mathrm{KL}}(Q_\theta\Vert Q_{\theta^{(j)}})/2}
 \leq\frac1{2\sqrt2}\leq\frac12.
 \label{eq:b:TV}
\end{equation}

Consider any measurable learner taking values in $\mathcal F$,
written $\widehat f=f_{\widehat v}$. For a fixed coordinate $j$, pair
each $\theta$ with $\theta_j=1$ with $\theta^{(j)}$, and write
$A_j=\{\widehat v_j=1\}$. The sum of the two testing errors satisfies
\begin{align*}
 Q_\theta(\widehat v_j\ne\theta_j)
 +Q_{\theta^{(j)}}(\widehat v_j\ne\theta_j^{(j)})
 &=Q_\theta(A_j^c)+Q_{\theta^{(j)}}(A_j)\\
 &\geq1-\operatorname{TV}(Q_\theta,Q_{\theta^{(j)}})\geq\frac12.
\end{align*}
There are $2^{d-1}$ such pairs. Averaging over all $2^d$ vertices gives
$2^{-d}\sum_\theta Q_\theta(\widehat v_j\ne\theta_j)\geq1/4$.
Sum this inequality over $j$ and use \eqref{eq:b:CVaR-hamming} to obtain
\[
 \frac1{2^d}\sum_\theta
 \mathbb E_{Q_\theta}\left[
 \rho_{\alpha,P_\theta}(\widehat f)-
 \inf_{f\in\mathcal F}\rho_{\alpha,P_\theta}(f)\right]
 \geq\frac{2\gamma}{d}\frac d4
 =\frac1{16}\sqrt{\frac{d}{n\pi_0\beta}}.
\]
The supremum over $\theta$ is at least this average, for every learner.
Taking the infimum over learners proves Theorem~\ref{thm:main:lower-bound}.
Randomized learners satisfy the same bound: append their independent
random seed to each experiment, which leaves the neighboring total
variation unchanged, and apply the identical testing argument.
\end{proof}

Only the expected $n\pi_0\beta$ revealed observations in the tail stratum
distinguish the vertices. This is an expected excess-CVaR lower bound for
$\mathcal F$-valued learners, already with fully identified risks at $\Gamma=1$.

